\chapter*{Zusammenfassung}
\addcontentsline{toc}{chapter}{Zusammenfassung}

\begin{otherlanguage}{ngerman}
Diese Arbeit beschreibt die Entwicklung und Evaluation des Prototyps \emph{Chatbot for Anatomy}. Das System verbindet dokumentenbasierte Fragebeantwortung mit einem getrennten Ablauf für den Export interaktiver 3D-Anatomie. Beide Ausgaben werden bewusst unterschiedlich behandelt, da sie unterschiedliche Nachweise benötigen. Hochgeladene Anatomie-PDFs werden mit MiniLM-Einbettungen in Chroma indexiert; ein CLIP-basierter Index kann zusätzlich Abbildungen liefern. Repository- und Zeitachsenbelege zeigen, dass die primäre englische Evaluation die hybrid-bevorzugte Produktionsroute unter \codepath{/rag/ask} nutzte: nach Commit \texttt{5089fcc} (12.~Juli~2026) konstruierte \codepath{MultimodalRAG} den \codepath{HybridRetriever} und fusionierte dichte, BM25- und Phrasen-Kandidaten mittels Reciprocal Rank Fusion, gefolgt von Deduplizierung, Reranking und Vorfilterung vor der Generierung. Ein früherer denser Snapshot (\texttt{cb041d5}, 27.~Juni~2026) mit einer festen Drei-Passagen-Grenze bleibt nur als historische Entwurfsstufe erhalten. Ein matched Vergleich denser versus hybrider Retrieval wurde nicht abgeschlossen; eine Überlegenheit des Hybridverfahrens wird daher nicht behauptet. Bei 3D-Anfragen löst das System eine Struktur über das Model Context Protocol und einen geometrisch validierten Z-Anatomy-Katalog auf, ruft Blender auf und liefert eine GLB-Datei, Annotationen sowie einen Three.js-Viewer. Ein Export wird nur dann als erfolgreich behandelt, wenn das Werkzeug ein gültiges Modellartefakt zurückgibt.

Die wichtigste englische Evaluation umfasste 55 Fragen; für 51 lagen verwertbare Antworten vor. Für 43 Nicht-Grenzfälle ergab ein automatisches Verfahren auf Basis der Einbettungsähnlichkeit Precision@3 von 0{,}86, Recall@3 von 0{,}58, nDCG@3 von 0{,}90 und MRR von 0{,}93. Diese Werte sind Ranking-Proxys und keine fachlichen Bewertungen von Relevanz, Antworttreue, Zitationsqualität oder anatomischer Richtigkeit. Die beobachteten Quellenzahlen hatten Mittelwert 5{,}16 und folgten fragentypabhängigen Grenzen von etwa vier bis acht statt der früheren festen Drei-Passagen-Regel. Dreizehn abgeschlossene Antworten wurden trotz einer dokumentierten Beschränkung auf die hochgeladenen Unterlagen als Nutzung von Weltwissen klassifiziert; dieses Feld wurde durch das Anfrageschema nicht durchgesetzt. Alle zehn bildbezogenen Anfragen lieferten Bildobjekte; manuelle Prüfungen während der Entwicklung zeigten jedoch unzuverlässige Relevanz, und eine systematische visuelle Evaluation wurde nicht abgeschlossen. Der MCP-Pfad führte zu einem nachvollziehbaren Erfolgskriterium, während die geplante vergleichende Laufzeitevaluation offenblieb. Der Beitrag der Arbeit liegt damit in einer dualen Architektur sowie in einer evidenzorientierten Darstellung der technischen Entscheidungen, beobachteten Fehler und noch ausstehenden Evaluationen.
\end{otherlanguage}

\cleardoublepage
